\answer{ex:02:1}{(а)~$\delta t_{\max}=0.64\ms$, $L=v\,\delta t_{\max}=3.2$~мм, шаг $v\cdot10\,\text{мкс}/2=25\um$, $\approx129$ детекторов. (б)~Срабатывает полоса шириной $v\,\tau\ln2=1.7$~мм, $\approx69$ детекторов; направление --- центр полосы.}
\answer{ex:02:2}{(а)~$-\tau\ln\frac{w_2}{\theta-w_1}\le\delta\le\tau\ln\frac{w_1}{\theta-w_2}$. (б)~$c=\frac\tau2\ln\frac{w_1(\theta-w_1)}{w_2(\theta-w_2)}$; окно $[-0.235,\,0.178]\ms$, $c=-28$~мкс: раньше должен прийти слабый вход. (в)~Направление смещается на $28$~мкс в сторону сильного уха; $c=0$ только при $\theta=w_1+w_2$, где окно нулевой ширины.}
\answer{ex:02:3}{(а)~$\sum_i\max(-z_i,0)$ удовлетворяет неравенству Грёнвалля с нулевым начальным значением и тождественно равно нулю. (в)~Нужны $s_i$ с $s_iw_{ij}s_j\ge0$; они существуют тогда и только тогда, когда циклы чётны. Взаимное торможение --- да (монотонна, устойчивых колебаний нет), петля \nt{ГЛУТ}--\nt{ГАМК} --- нет (может колебаться).}
\answer{ex:02:4}{(б)~Шесть нейронов $g_k=[x+y+z\ge k]$ и $\bar g_k=[\bar x+\bar y+\bar z\ge4-k]$, $k=1,2,3$; $C=g_2$, $\bar C=\bar g_2$, $S=[g_1+\bar g_2+g_3\ge2]$, $\bar S=[\bar g_1+g_2+\bar g_3\ge2]$: восемь нейронов. (г)~$12$ нейронов.}
\answer{ex:02:5}{$r_0=0.00197$; $T_{\min}(0.6)=0.718\,\tau$; граница $\tau\ln[2(1-r_0)]=0.691\,\tau$; $I_c=0.225$, $T_{\min}\propto(I-I_c)^{-1/2}$.}
\answer{ex:02:6}{(а)~$500$ звеньев, $5\cdot10^4$ нейронов и импульсов. (б)~$4.5\ms$ ($0.45\,\%$) за $1\,\text{с}$, $1.4\ms$ ($1.4\,\%$) за $100\ms$: относительная ошибка $\propto T^{-1/2}$. (в)~Общий для всех звеньев случайный множитель задержки с разбросом $5\,\%$.}
\answer{ex:02:7}{(а)~Освежать раз в $\tau_F\ln2=0.69\,\text{с}$: $4.3$ импульса в секунду на нейрон против $20$; $4.3\cdot10^3$ импульсов ($4\cdot10^{-7}$~Дж) против $2\cdot10^4$ ($2\cdot10^{-6}$~Дж). (б)~Триггер теряет бит; конденсатор сохраняет его, если глушение началось не позже чем через $0.49\,\text{с}$ после освежения (вероятность $0.71$), а наверняка --- при освежении раз в $0.49\,\text{с}$, $6.1$ импульса в секунду.}
